\section{研究背景与问题定义}
\label{sec:background}

\subsection{从工业开放问题到可计算研究对象}

众包离线定位利用附近设备接收到的无线广播，为暂时无法联网的目标提供位置估计。课程参考材料《HarmonyOS难题发布》第七期第1页将这一任务表述为“基于非线性优化的众包离线高精度位置聚合难题”\cite{huaweiHarmonyChallenge7}：目标持续发送蓝牙广播，协助设备上报自身定位与接收信号强度；系统需要从时间、定位精度和可信程度不同的报告中估计目标当前位置。困难并不只来自测距噪声。报告者坐标本身存在误差，历史报告可能对应不同的目标位置，少量严重错误上报还可能破坏聚类和非线性拟合。

原材料给出的静态模型为
\begin{equation}
  \widehat x=\operatorname*{arg\,min}_{x\in\mathbb R^2}
  \sum_{i=1}^{N}\omega_i\bigl(f(r_i)-\|x-l_i\|_2\bigr)^2,
  \label{eq:original-problem}
\end{equation}
其中 $l_i$ 为报告者位置，$r_i$ 为接收信号强度，$f(r_i)$ 为反演距离，$\omega_i$ 由定位时间、方式、精度和置信度等信息决定。材料提出，超过10个协助设备时CEP90应在30米以内，并要求在标准场景不劣化的条件下将移动目标检测精度提升20\%。后一指标没有给出具体统计定义；材料也未提供原系统实现和真实验收数据。因此，不能用某个合成基准上的平均误差下降替代这项工业要求。

本文将研究对象限定为\emph{受污染、异步、二维距离观测下的当前位置估计及其可信度评估}。报告者可移动，但每条上报包含该次观测时的位置；目标在一个短时间窗口内以静态或恒速模型近似。输入是上报坐标、时间、距离及标定尺度，输出是当前位置、求解诊断和可选的预测区域。无线协议、实际硬件标定、楼层判断、地图约束和大规模产品部署不在本次实现范围内。这种限定保留了原题的非线性估计、运动失配、污染和几何退化等核心数学困难，同时允许在本地CPU上重复实验。

\subsection{相关研究及本文的技术位置}

接收信号强度与距离的关系通常在对数域建模。Patwari等讨论了RSS与到达时间观测下的位置估计、统计模型和误差界\cite{patwari2003relative}。当分贝域噪声近似高斯时，反演距离具有乘性误差，因此距离域的平方残差与对数残差对应不同的统计假设。报告者坐标误差、非视距传播和设备标定差异还会进一步改变有效噪声。本项目借助这一背景构建可控模拟，并不将其视为对真实BLE环境的充分刻画。

稳健估计为少量错误观测提供了已有工具。Soares与Gomes将Huber损失用于受污染的测距定位，并研究相应非凸问题的凸下界和分布式算法\cite{soares2016badapples}。Fischler与Bolles提出的RANSAC则通过子集拟合与全体一致性评价寻找不受大量粗差支配的模型\cite{fischler1981random}。本文采用的非线性稳健损失、子集初始化和截尾重拟合受这些思想启发，但数值实现没有继承相关文献中另一种算法的全局最优或收敛保证。

选择静态还是运动模型是一项预测问题。交叉验证提供了区分训练拟合和观测外预测的基本方法\cite{stone1974crossvalidatory}；它本身不会消除小样本、相依观测和模型失配。定位唯一性还受到几何信息限制，局部优化是否收敛与观测是否能区分两个不同位置是两回事\cite{eren2004rigidity}。对于不确定性，split conformal能够在交换性和方法冻结等条件下提供边际预测覆盖\cite{lei2018distributionfree}，但不能直接保证任一具体困难样本或任意拒绝后子集的条件覆盖。

由此，本文并不把某个现成损失函数称为新的数学发现。可检验的研究内容是：这些组件在同一个异步定位问题中如何相互作用，哪些失败可由精确反例解释，以及工具使用型科研智能体能否根据真实证据作出有依据的修订。

\subsection{先导实验暴露的失败}

第一版原型采用固定候选库中的经验规划，尚未调用大语言模型。其独立合成测试结果见表\ref{tab:pilot-failure}。这些数字属于历史先导实验，不属于本文后续新版系统的最终测试。

\begin{table}[tbp]
  \centering
  \caption{第一版先导实验的历史结果：160个合成测试案例。误差越低越好。}
  \label{tab:pilot-failure}
  \begin{tabular}{lrr}
    \toprule
    方法 & 平均误差（m） & CEP90（m）\\
    \midrule
    静态WLS & 6.4840 & 14.3271\\
    经验规划所选配置 & 4.3632 & 12.4707\\
    同预算随机搜索所选配置 & 3.6839 & 10.2265\\
    固定恒速稳健模型 & 4.1668 & 9.5568\\
    \bottomrule
  \end{tabular}
\end{table}

经验规划虽然相对静态WLS改善总体均值，却没有优于随机搜索或固定恒速稳健模型。在移动异常场景中，启发式模型评分将40例中的37例选择为静态，该场景均值为11.2833米，略高于WLS的11.2000米；静态正常场景也从0.8743米退化到0.9113米。这表明，总体均值下降既不能证明搜索策略更有效，也不意味着实现了原题的标准场景非劣要求。

这些观察引出三个待检验解释：由观测距离自身构造尺度可能改变误差权重；污染下的训练拟合评分可能误判运动；有限开发案例上的选择可能发生泛化失败。它们不是已经确定的因果结论。只有保持其余组件不变的消融，以及独立案例上的验证，才能区分各项解释。因为这些历史测试结果已被研究者看到，它们可以作为诊断材料，不能再次被称为新版方案完全未见的最终测试。

\subsection{研究问题与贡献边界}

本文围绕以下相互关联的问题组织研究。
\begin{enumerate}
  \item 对数残差及其正确中心化，能否改善乘性误差下的拟合，并减少观测依赖权重导致的系统偏差？
  \item 在相同求解组件下，观测交叉验证能否比训练残差驱动的门控更好地权衡静态稳定性与运动解释能力？
  \item 局部几何诊断、近似预测尺度与独立校准，能够支持什么范围的可信度结论；哪些歧义无法仅靠数值优化消除？
  \item 一个实际调用语言模型、实验与核验工具的科研控制器，是否形成可复盘的假设、行动、观察与修订链条？
\end{enumerate}

数学贡献采用明确假设下的项目内推导：标量相对平方风险的偏差反例、恒速镜像歧义及不可区分性下界、包含曲率和邻域条件的局部误差界。它们用于解释与约束实验，不作理论首创声明。系统贡献是独立实现的工具调用运行时、可组合定位算子、证据与状态记录，以及真实推理和日志回放的区分。

架构参考来自Creative-Intelligence所强调的研究、验证和持久状态分离思想\footnote{参考项目：\texttt{github.com/DeepMathLLM/Creative-Intelligence}，查阅日期为2026年9月25日。}。本项目不是该仓库的fork，没有复制其源码，也不以运行Moonshine作为前提。本文对智能体实际行为的描述以保存的模型请求、响应和工具证据为依据；模型给出的评论既不是独立实验复现，也不是数学证明。
